\chapter{Introduction}
\label{sec:intro}

\section{Motivation}
\label{sec:intro:mo}

Business processes represent core operational assets of modern organizations, defining standardized workflows, responsibilities, and execution logic for daily operations. However, a large proportion of enterprise process knowledge remains trapped in unstructured document formats, such as Word, PDF, plain text, and internal manuals. Manually converting such unstructured textual content into formal, executable process models is labor-intensive, time-consuming, and prone to human error, reducing the efficiency of process design and digitization \cite{daclin2024gai4bm}.

In the era of generative artificial intelligence (AI), large language models (LLMs) have emerged as powerful enablers for automated process modeling, demonstrating strong capabilities in understanding natural language, extracting logical structures, and generating standardized process notations. Nevertheless, LLMs cannot fully replace human experts, especially in domain-specific scenarios requiring professional knowledge, contextual judgment, and structural decision-making \cite{klievtsova2024process}. Therefore, hybrid human-LLM process modeling has become a promising paradigm that combines human domain expertise with AI efficiency. 

Despite rapid advancements, existing LLM-assisted process modeling solutions suffer from critical limitations in document-driven scenarios. Most tools adopt a “single-view, single-model” paradigm, which lacks scalability for complex systems involving multiple processes and subprocesses. To overcome these gaps, this thesis proposes an integrated tool, Document-Based Process Modeler. The tool enables a seamless, efficient, and user-centric modeling workflow with three core innovations: unified workspace management for correlated documents and models; direct text selection from uploaded documents as LLM input; and side-by-side visualization to establish explicit visual link between source text and generated process model to make a more intuitive and verifiable modeling experience.

\section{Research Questions}

\label{sec:intro:rq}

This thesis addresses the following research questions:

RQ1 How do existing tools support AI-assisted process modeling multiple documents as input?

RQ2: What functionalities should a system provide to support AI-assisted process modeling from multiple documents, and how should these functionalities be designed and integrated in a seamless and user-friendly manner?

RQ3 How do user perceive the benefits and limitation of the system?

\section{Contribution}
\label{sec:intro:con}
The contribution of this thesis consists of the following aspects:

Firstly, a systematic review of the existing tools and approaches for AI-assisted process modeling is conducted, which provides insights into the current state of the art and guides the design of the proposed system.

Secondly, an integrated web-based application\footnote{\url{https://github.com/ga94hor/doc-based-process-modeler}, accessed 2026-04-15} is designed and developed to support process modeling from documents using LLMs. The system offers a more seamless, efficient, and user-friendly workflow by integrating multiple functionalities, such as document management, model generation, and traceability between text and model.

Thirdly, the proposed system is evaluated through a user study with process modeling experts. The evaluation results provide insights into the system’s benefits and limitations from the users’ perspective, which can inform future improvements and further research in this area.

\section{Methodology}
\label{sec:intro:meth}
The research process of this thesis follows the three-cycle view of Design Science Research (DSR) proposed by Hevner \cite{hevner2007}. This framework consists of three interconnected cycles: the relevance cycle, the design cycle, and the rigor cycle, which link the problem environment, artifact development, and the knowledge base. In this work, these cycles are applied as follows.

The relevance cycle is initiated by problems in current process modeling tools, especially the lack of support for document-based modeling. These problems are continuously discussed with modeling experts throughout the development process to identify practical needs and define system requirements.

The design cycle focuses on designing, implementing, and improving the artifact (RQ2), the Document-Based Process Modeler. The system is developed and evaluated iteratively. Feedback from the evaluation (RQ3) is used to refine and improve the design.

The rigor cycle is supported by continuously reviewing state-of-the-art AI-assisted process modeling tools in the literature review (RQ1), ensuring that the design is grounded in prior work and contributes to the research field.

\section{Evaluation}
\label{sec:intro:ev}
The evaluation of the proposed system is done with 4 experts in qualitative way. It consisted of task-based activities and a follow-up questionnaire to collect user feedback on the tool's usability and perceived usefulness.

\section{Structure}
\label{sec:intro:struct}
The remainder of this thesis is structured as follows: Chapter~\ref{sec:rel} reviews the related work in AI-assisted process modeling tools. Chapter~\ref{sec:solution} presents the design solution of the tool, including its requirements, architecture, and functionalities. Chapter~\ref{sec:implementation} describes the technical implementation of the proposed system. Chapter~\ref{sec:evaluation} details the evaluation, including the methodology and results. Chapter~\ref{sec:discussion} discusses the implications of the findings and potential future work. Finally, Chapter~\ref{sec:conclusion} summarizes the thesis.
